% Q2V2: Parallel RLC Step Response, Overdamped — Exam diagram
% Current source Is switches on at t=0, parallel R||L||C
\documentclass[border=10pt]{standalone}
\usepackage[american voltages, american currents]{circuitikz}
\usepackage{amsmath}
\renewcommand{\familydefault}{\sfdefault}

\begin{document}
\begin{circuitikz}[line width=0.8pt, every node/.style={font=\sffamily}]

% === Current source with closing switch ===
\draw (0,0) to[I, l=$I_s$, i=$i_s$] (0,3.5);
\draw (0,3.5) to[closing switch] (0,5);

% Switch labels (separate nodes to avoid overlap)
\node[left, font=\sffamily\itshape\small] at (-0.5,4.25) {SW1};
\node[left, font=\sffamily\footnotesize] at (-0.5,4.7) {$t{=}0$};
\node[left, font=\sffamily\footnotesize] at (-0.5,3.8) {(closes)};

% === Top wire ===
\draw (0,5) -- (8,5);
% === Bottom wire ===
\draw (0,0) -- (8,0);

% === Parallel branches ===
\draw (2.5,5) to[R, l=$R$, i>^=$i_R$] (2.5,0);
\draw (5,5) to[L, l=$L$, i>^=$i_L$] (5,0);
\draw (8,5) to[C, l=$C$, i>^=$i_C$] (8,0);

% === Voltage across combination (offset further right for clearance from C label) ===
\draw (10,5) to[open, v^=$v(t)$] (10,0);
\draw (8,5) -- (10,5);
\draw (8,0) -- (10,0);

% === Junction dots ===
\fill (0,0) circle (2pt); \fill (0,5) circle (2pt);
\fill (2.5,0) circle (2pt); \fill (2.5,5) circle (2pt);
\fill (5,0) circle (2pt); \fill (5,5) circle (2pt);
\fill (8,0) circle (2pt); \fill (8,5) circle (2pt);

\end{circuitikz}
\end{document}
